\icode{E-ME-SI-L3-261017-001}
全球最大水平臂上回转自升塔式起重机的开发和应用，意味着中国桥梁及铁路施工装备进一步迈向世界前列。该起重机某次从\(t = 0\)时刻由静止开始提升质量为\(m\)的物体，其所受合外力随时间变化的图像如图所示，\(t_1\)到\(t_2\)内起重机的功率为额定功率，不计物体受到的空气阻力，重力加速度为\(g\)，下列说法正确的是\paren[]

\begin{tikzpicture}[>=stealth,font=\scriptsize]
	\draw[->,blue] (-.5,0)--(3.5,0)node[below]{\(t\)/\qty{}{s}};
	\draw[->,blue] (0,-.5)--(0,3)node[right]{\(F\)/\qty{}{N}};
	\node[blue] at(-.2,-.2){\(O\)};
    \draw[dashed,cyan] (1,2)--(1,0) (2.2,{2*exp(-2.2)/exp(-1)})--(2.2,0);
	\draw[magenta] (0,2)--(1,2);
	\draw[domain=1:3,smooth,magenta] plot(\x,{2*exp(-\x)/exp(-1)});
	\node[blue] at(0,2)[left]{\(F_0\)};
	\node[blue] at(1,0)[below]{\(t_1\)};
	\node[blue] at(2.2,0)[below]{\(t_2\)};
\end{tikzpicture}
\begin{choices}
	\item 物体匀加速阶段的加速度为\(\dfrac{F_0 - mg}{m}\)
	\item 0到\(t_1\)和\(t_1\)到\(t_2\)时间内牵引力做的功之比为\(\dfrac{t_1}{2(t_2 - t_1)}\)
	\item \(t_2\)时刻物体正在减速上升
	\item 0到\(t_1\)阶段牵引力所做的功为\(\dfrac{F_0^2 t_1^2}{2m}\)
\end{choices}
\begin{answer}
    B\qquad 提示：匀加速阶段功率从0均匀增加到额定功率\(P_0\)，\(0\sim t_1\)的平均功率为\(\overline{P} =\frac{P_0}{2}\)
\end{answer}